\begin{frame}{그림으로 읽기: 고도로 벌리면 평면이 자른다}
  \begin{columns}[c]
    \begin{column}{0.60\textwidth}
      \centering
      \includegraphics[width=\linewidth]{ch05_3_circles_3d_lift.png}
    \end{column}
    \begin{column}{0.37\textwidth}
      \small
      \begin{itemize}
        \item 좌: 동심원 $\to$ $z=x_1^2+x_2^2$ 축으로 들어올림.
        \item 우: 두 원환체, \textbf{평면 $z=0.5$} 하나로 자름.
        \item 평면의 2차원 투영 $=$ \textbf{원형} 경계.
      \end{itemize}
      \vskip0.5em
      \begin{block}{일반 원리}
        \textbf{곡면으로만} 나뉘는 관계 $\to$ 곡면의 고도를 새
        특징으로 추가 $\to$ \textbf{평면으로 나뉘는} 관계.
      \end{block}
      {\footnotesize XOR도 동일: $(x_1,x_2,x_1x_2)$로 올리면
      평면 분리 문제.}
    \end{column}
  \end{columns}
\end{frame}
